% Nouveau EGA V (cote 25): archive cote 25, site page 48 (PDF page 49).
% typescript. Draft compared: none.
% AI-transcribed by Claude Opus 5.5 (claude-opus-5-5, adaptive thinking, effort high), Message Batches API, batch ega25b (msgbatch_01A7n5YuXubUPNK1zsepi8RG), 2026-11-13; original AI transcription, subsequently checked against both supplied scan parts; see PAGES_CHECKED.tsv and CHECK_DECISIONS.tsv.
\scanpage{25/48}
% NOTE: Matching loop-headed reference marks carry the top grâce note to the proof after codimension2. Insertion hooks are not mathematical parentheses or slashes. The lower soit contenue dans X prime is legible but its attachment and the force of its small enclosure remain unsettled.

\mspage{\hand{\unsure{44|45}}}




que $\phi_{\hand{\xi}}$ soit $X_{k(\hand{\xi})}$-régulière et $Y_{\hand{\xi}}$ satisfasse la condition $R_k$ géométri-%

que. Alors:

a) Bour que $Z_{k-1}$ soit de codimension $\hand{\geq} 2$ dans $\overset{\hand{\vee}}{P}$ , il  s que
 
toute composante irréductible $X'_i$ de $X'$ de dim $\leq k$ soit \uline{réglée} pour $f$. 

b) Pour que $Z_k$ soit de codimension $\hand{\geq} 2$ dans $\overset{\hand{\vee}}{P}$ , il  s que

toute composante irréductible 
$X_i$ de dim $\leq k+1$ de $X$ soit formée de pts lisses de $X$ et soit \uline{réglée} pour $f$.



 
soit contenue dans $X'$

En effet, pour tout $\hand{\xi}$, l'ensemble géométriquement singulier

de $Y_{\hand{\xi}}$  (NB on se borne aux $\hand{\xi}$ tels que $\phi_{\hand{\xi}}$ soit $X_{k(\hand{\xi})}$-

régulier, ce qui est inoffensif vu 8.2.) est la réunion de

sing$(Y'_{\hand{\xi}}$ ) et de l'image inverse $T_{\hand{\xi}}$ de $T$ dans $Y_{\hand{\xi}}$ , donc la codimension
\placedmargin{lieu singulier géométrique}
de cet ensemble singulier dans $Y_{\hand{\xi}}$ est égale à l'inf de codim (sing$(Y'_{\hand{\xi}})$,$Y'_{\hand{\xi}}$ )

et de codim $(T_{\hand{\xi}},Y_{\hand{\xi}})$. Bornons nous aux $\hand{\xi}$ tels que \hand{)}
\placedmargin{$*$} sing$(Y'_\xi)$ soit fini
(ce qui est inoffens\unsure{i|f}f, car  cela condu\unsure{i|t}t à se mettre dans le

complémentaire d'un ensemble de \unsure{codimension} $\hand{\geq} 2$ grâce au fait que $\underline{X'}^{\text{sing}}$ est de dim $r-1$, où $r=\dim \check{P}$ !), Les points singulier

géométriques de $Y'_{\hand{\xi}}$ sont alors isolés, La conclusion résulte facilement de

là , et de 8.4. 

Combinant 8.13.  8.5. et 8.6., on trouve à la façon habituelle:

\uline{Corollaire} 8.14. On suppose $f:X\hand{\rightarrow} P$ non ramifié et $X$ sans pts isolés.

\placedmargin{n} a) Supposons $X$ séparable sur $k$. Pour que l'ensemble  des $\hand{\xi \in}$ $\overset{\hand{\vee}}{P}$

tels que $\phi_{\hand{\xi}}$ soit $X_{k(\hand{\xi})}$ -régulier et $Y_{\hand{\xi}}$ soit séparable ait un complé-%

mentaire de codimension $\hand{\geq} 2$, il f et s qu'

\placedmargin{\hand{48}}
